\QHeading{第1题　口算。}
\par\noindent


\begin{oralgrid}
\OralItem{1}{32\times3}{96}
\OralItem{2}{450\div50}{9}
\OralItem{3}{60\times15}{900}
\OralItem{4}{840\div40}{21}
\OralItem{5}{17\times5}{85}
\OralItem{6}{240\div6}{40}
\OralItem{7}{35\times2}{70}
\OralItem{8}{22\times4}{88}
\OralItem{9}{15\times60}{900}
\OralItem{10}{75\div5}{15}
\OralItem{11}{70\times30}{2100}
\OralItem{12}{60\div4}{15}\end{oralgrid}
\par

\QHeading{口算思路补充（对应上表题号）}

\Answer{上面列出得数，下面按小题给出不同计算思路。计数单位法适用于小数，整数题可比较凑整、拆分或商不变等方法。}
\MethodTitle{第1小题：$32\times 3$}

\MethodBlock{口算思路一：凑整后补偿}{把接近整十、整百的因数凑整，再补回或减去多算的部分。}{\begin{aligned}[t]&\phantom{={}}32\times 3\\
&=3\times (30+2)\\
&=3\times 30+3\times 2\\
&=90+6\\
&=96\end{aligned}}
\MethodTitle{第2小题：$450\div 50$}

\MethodBlock{口算思路一：同时去掉相同个数的0}{}{\begin{aligned}[t]&\phantom{={}}450\div 50\\
&=(450\div 10)\div (50\div 10)\\
&=45\div 5\\
&=9\end{aligned}}

\MethodBlock{口算思路二：利用商不变的性质}{被除数和除数同时乘同一个非零数，商不变。}{\begin{aligned}[t]&\phantom{={}}450\div 50\\
&=(450\times 2)\div (50\times 2)\\
&=900\div 100\\
&=9\end{aligned}}
\MethodTitle{第3小题：$60\times 15$}

\MethodBlock{口算思路一：按数位拆分计算}{把一个因数拆成整十、整百数与余下部分之和。}{\begin{aligned}[t]&\phantom{={}}60\times 15\\
&=60\times (10+5)\\
&=60\times 10+60\times 5\\
&=600+300\\
&=900\end{aligned}}
\MethodTitle{第4小题：$840\div 40$}

\MethodBlock{口算思路一：同时去掉相同个数的0}{}{\begin{aligned}[t]&\phantom{={}}840\div 40\\
&=(840\div 10)\div (40\div 10)\\
&=84\div 4\\
&=21\end{aligned}}

\MethodBlock{口算思路二：拆分被除数}{}{\begin{aligned}[t]&\phantom{={}}840\div 40\\
&=(800+40)\div 40\\
&=800\div 40+40\div 40\\
&=20+1\\
&=21\end{aligned}}
\MethodTitle{第5小题：$17\times 5$}

\MethodBlock{口算思路一：按数位拆分计算}{把一个因数拆成整十、整百数与余下部分之和。}{\begin{aligned}[t]&\phantom{={}}17\times 5\\
&=5\times (10+7)\\
&=5\times 10+5\times 7\\
&=50+35\\
&=85\end{aligned}}
\MethodTitle{第6小题：$240\div 6$}

\MethodBlock{口算思路一：利用商不变的性质}{被除数和除数同时乘同一个非零数，商不变。}{\begin{aligned}[t]&\phantom{={}}240\div 6\\
&=(240\times 5)\div (6\times 5)\\
&=1200\div 30\\
&=40\end{aligned}}

\MethodBlock{口算思路二：拆分除数，连续相除}{}{\begin{aligned}[t]&\phantom{={}}240\div 6\\
&=240\div (2\times 3)\\
&=240\div 2\div 3\\
&=120\div 3\\
&=40\end{aligned}}
\MethodTitle{第7小题：$35\times 2$}

\MethodBlock{口算思路一：按数位拆分计算}{把一个因数拆成整十、整百数与余下部分之和。}{\begin{aligned}[t]&\phantom{={}}35\times 2\\
&=2\times (30+5)\\
&=2\times 30+2\times 5\\
&=60+10\\
&=70\end{aligned}}
\MethodTitle{第8小题：$22\times 4$}

\MethodBlock{口算思路一：凑整后补偿}{把接近整十、整百的因数凑整，再补回或减去多算的部分。}{\begin{aligned}[t]&\phantom{={}}22\times 4\\
&=4\times (20+2)\\
&=4\times 20+4\times 2\\
&=80+8\\
&=88\end{aligned}}
\MethodTitle{第9小题：$15\times 60$}

\MethodBlock{口算思路一：按数位拆分计算}{把一个因数拆成整十、整百数与余下部分之和。}{\begin{aligned}[t]&\phantom{={}}15\times 60\\
&=60\times (10+5)\\
&=60\times 10+60\times 5\\
&=600+300\\
&=900\end{aligned}}
\MethodTitle{第10小题：$75\div 5$}

\MethodBlock{口算思路一：拆分被除数}{}{\begin{aligned}[t]&\phantom{={}}75\div 5\\
&=(50+25)\div 5\\
&=50\div 5+25\div 5\\
&=10+5\\
&=15\end{aligned}}

\MethodBlock{口算思路二：利用商不变的性质}{被除数和除数同时乘同一个非零数，商不变。}{\begin{aligned}[t]&\phantom{={}}75\div 5\\
&=(75\times 2)\div (5\times 2)\\
&=150\div 10\\
&=15\end{aligned}}
\MethodTitle{第12小题：$60\div 4$}

\MethodBlock{口算思路一：拆分被除数}{}{\begin{aligned}[t]&\phantom{={}}60\div 4\\
&=(40+20)\div 4\\
&=40\div 4+20\div 4\\
&=10+5\\
&=15\end{aligned}}

\MethodBlock{口算思路二：利用商不变的性质}{被除数和除数同时乘同一个非零数，商不变。}{\begin{aligned}[t]&\phantom{={}}60\div 4\\
&=(60\times 5)\div (4\times 5)\\
&=300\div 20\\
&=15\end{aligned}}


\par\noindent

\QHeading{第2题　脱式计算。}

\CalcBlock[方法一]{1}{\begin{aligned}[t]&\phantom{={}}\color{black}[36-(74-58)]\times35\\&\color{answerred}=[36-16]\times35\\&\color{answerred}=20\times35\\&\color{answerred}=700\end{aligned}}
\par
\MethodBlock{方法二：乘法分配律}{分别求出两部分的积，再合起来。}{\begin{aligned}[t]&\phantom{={}}[36-(74-58)]\times 35\\
&=36\times 35-(74-58)\times 35\\
&=36\times 35-16\times 35\\
&=1260-560\\
&=700\end{aligned}}


\CalcBlock{2}{\begin{aligned}[t]&\phantom{={}}\color{black}96\div[960\div(215-135)]\\&\color{answerred}=96\div[960\div80]\\&\color{answerred}=96\div12\\&\color{answerred}=8\end{aligned}}
\CalcBlock{3}{\begin{aligned}[t]&\phantom{={}}\color{black}(126+121)\div19\times72\\&\color{answerred}=247\div19\times72\\&\color{answerred}=13\times72\\&\color{answerred}=936\end{aligned}}
\CalcBlock{4}{\begin{aligned}[t]&\phantom{={}}\color{black}(85\times23-1707)\div31\\&\color{answerred}=(1955-1707)\div31\\&\color{answerred}=248\div31\\&\color{answerred}=8\end{aligned}}
\CalcBlock{5}{\begin{aligned}[t]&\phantom{={}}\color{black}[99+(305-248)]\div52\\&\color{answerred}=[99+57]\div52\\&\color{answerred}=156\div52\\&\color{answerred}=3\end{aligned}}
\CalcBlock{6}{\begin{aligned}[t]&\phantom{={}}\color{black}85\times[(233-89)\div12]\\&\color{answerred}=85\times[144\div12]\\&\color{answerred}=85\times12\\&\color{answerred}=1020\end{aligned}}
\par\noindent

\QHeading{第3题　解决问题。}

\Problem{①文具店购进彩色卡纸，每包的单价为18元。王老师打算一次性购买32包，带600元够吗？}
\par


\MethodTitle{方法一}


\MethodTitle{方法二：用整十单价作上界判断}
\Answer{每包18元比20元少2元。按20元买32包会超过600元40元，但实际少花64元，所以够。}
\Answer{\textbf{按每包20元多算的费用}\par \Formula{\begin{aligned}[t]&\phantom{={}}(20-18)\times 32\\
&=2\times 32\\
&=64\text{（元）}\end{aligned}}}
\Answer{\textbf{按20元计算超出所带的钱}\par \Formula{\begin{aligned}[t]&\phantom{={}}20\times 32-600\\
&=640-600\\
&=40\text{（元）}\end{aligned}}}
\Answer{\textbf{实际还剩}\par \Formula{\begin{aligned}[t]&\phantom{={}}64-40\\
&=24\text{（元）}\end{aligned}}}
\Answer{答：带600元够，还剩24元。}
\Problem{②商店里每个文具盒15元，每把直尺3元。四(2)班购买24个文具盒和40把直尺，一共需要花多少钱？}

\MethodTitle{方法一}
\Answer{\Formula{\begin{aligned}[t]&\phantom{={}}15\times24+3\times40\\&=360+120\\&=480\text{（元）}\end{aligned}}\par 答：一共需要花480元。}

\MethodTitle{方法二：把文具盒费用折算成3元一份}
\Answer{15元是3元的5倍，24个文具盒的费用相当于120份3元，再加40把直尺。}
\Answer{\Formula{\begin{aligned}[t]&\phantom{={}}3\times (5\times 24+40)\\
&=3\times (120+40)\\
&=3\times 160\\
&=480\text{（元）}\end{aligned}}}
\Answer{答：一共需要480元。}


\Answer{
\QHeading{计算过程：第3题①：购买彩色卡纸}
购买32包需要：
\Formula{\begin{aligned}[t]&\phantom{={}}18\times32\\&=18\times(30+2)\\&=18\times30+18\times2\\&=540+36\\&=576\text{（元）}\text{。}\end{aligned}}
\(576<600\)，所以带的钱够，还剩\(600-576=24\)元。

答：够，还剩24元。


}
